\chapter{Sources}

\section*{The objects}

Josephine Nivison Hopper and Edward Hopper, \emph{Artist's ledger --- Book I}, 1913--1963. Whitney Museum of American Art, 96.208. Gift of Lloyd Goodrich. 117 digitised sheets.

Josephine Nivison Hopper and Edward Hopper, \emph{Artist's ledger --- Book II}, 1907--1962. Whitney Museum of American Art, 96.209. Gift of Lloyd Goodrich. 72 sheets.

Josephine Nivison Hopper and Edward Hopper, \emph{Artist's ledger --- Book III}, 1924--1967. Whitney Museum of American Art, 96.210. Gift of Lloyd Goodrich. 90 sheets.

\emph{Artist's ledger --- Book IV}, 1913--1967. Whitney Museum of American Art, 96.211. Purchase. 161 sheets. Hand unattributed.

Josephine Nivison Hopper and Edward Hopper, \emph{Artist's ledger --- Book V}, 1953--1963. Whitney Museum of American Art, 96.212. Gift of Lloyd Goodrich. 16 sheets.

Josephine Nivison Hopper, \emph{Artist's ledger --- Dealers/Etchings}, 1921--1951 (contents 1913--1961). Whitney Museum of American Art, 96.213. Purchase, with funds from an anonymous donor. 48 sheets.

Josephine Nivison Hopper, \emph{Garrulities + Grouch} (1924--1950, contents to 1959), 42 sheets; \emph{Re: 3 Wash Sq} (1947), 51 sheets; \emph{Battle of Wash Sq.} (1947, contents to 1955), EJHA.1597, 8 sheets; black two-ring notebook (circa 1952--1961, contents 1937--1964), 68 sheets. Whitney Museum of American Art, Sanborn Hopper Archive, Series IV, Subseries A.

Josephine Nivison Hopper, diaries 1933--1956, 22 volumes. Provincetown Art Association and Museum, gift of 2016. Cited through Madeleine Larson's typescript, Parts I--V, 389 pages, by part and page, as the edition's timeline cites them.

\section*{The reading}

\emph{Hopper Ledgers}, \url{https://hopper.commutator.io}. Transcriptions of the six ledgers (46 batches, 543 sheets) and the four notebooks, first machine pass by Claude Opus 5 and Claude Fable 5.1, 5 to 13 September 2026, unchecked against the photographs by a person except where a leaf's note says otherwise. This critical note was itself drafted by Claude Fable 5.1 on 13 September 2026, under the direction of Michel Hua, from those transcriptions and the derived files; see the front note and Chapter 8. The \texttt{.tex} under \texttt{transcripts/} is the source of record; the tables in the appendix are read from \texttt{src/content/}. Method: \texttt{/method/}; Book IV's reader: \texttt{docs/book-iv-method.md}; the glossary of forms: \texttt{/schema/}.

Every citation in this draft names the leaf as numbered on the paper and the Whitney's ResourceSpace reference for the photograph. A sheet is at \texttt{hopper.commutator.io/} followed by the volume, a hash, and the volume, batch and ref, and its photograph is fetched from the museum as it is looked at; nothing of the archive is stored by the edition.

\section*{Prior readers}

Gail Levin, \emph{Edward Hopper: An Intimate Biography} (1995), written from the diaries; \emph{Edward Hopper: A Catalogue Raisonné} (Whitney Museum of American Art and W. W. Norton, 1995), written from the record books.

Deborah Lyons, \emph{Edward Hopper: A Journal of His Work} (Whitney Museum of American Art and W. W. Norton, 1997), the record books published.

Madeleine Larson, typed transcripts of Josephine Hopper's diaries, Parts I--V, Provincetown Art Association and Museum.

\section*{Museum records cited}

The Art Institute of Chicago, 1942.51 (Nighthawks), 1939.2082 (The Evening Wind), 1923.305 (East Side Interior); The Metropolitan Museum of Art, 25.31.1--15 (fifteen etchings, Harris Brisbane Dick Fund, 1925), the Hearn Fund purchase of 1931 (Tables for Ladies); The Museum of Modern Art, 3.1930 (House by the Railroad); Yale University Art Gallery, 1961.18.29 (Rooms by the Sea); Des Moines Art Center, 1958.2 (Automat). Each was retrieved by the edition's \texttt{scripts/works.mjs} from the museum's own API and checked against the artist field; the per-work notes in \texttt{work-notes.json} carry the retrieval date.
